% TOP_PROBLEM: {"id": 20001087, "title": "A two-sided transfer principle for the least doubling of 3-AP-free sets", "category_id": 2, "category": {"id": 2, "name": "combinatorics", "display_name": "Combinatorics", "description": "Counting problems, graph theory, discrete structures.", "slug": "combinatorics", "order_index": 2, "created_at": "2026-07-31T15:26:25.671Z"}, "status": "partially_solved", "problem_number": "AIM-COMBINATORICS-0212", "set_id": 14}
% FIRST_POSTED: 2026-10-01
% ENTRY_KIND: statement_only
% UnsolvedMath ID 20001087; source catalogue code AIM-COMBINATORICS-0212
% !TeX program = lualatex
% Definition and sources only; this file is not an LLM solution attempt.
\documentclass[11pt,a4paper]{article}
\usepackage[margin=25mm]{geometry}
\usepackage{fontspec}
\setmainfont{lmroman10-regular.otf}[BoldFont=lmroman10-bold.otf,
 ItalicFont=lmroman10-italic.otf,BoldItalicFont=lmroman10-bolditalic.otf]
\usepackage{amsmath,amssymb,mathtools}
\usepackage{unicode-math}
\setmathfont{latinmodern-math.otf}
\AtBeginDocument{\let\setminus\smallsetminus}
\usepackage{xurl,hyperref}
\hypersetup{colorlinks=true,urlcolor=blue}
\setcounter{secnumdepth}{0}
\setlength{\parindent}{0pt}
\setlength{\parskip}{5pt}
\setlength{\emergencystretch}{3em}
\begin{document}
\section*{A two-sided transfer principle for the least doubling of 3-AP-free sets}\label{20001087}
{\small UnsolvedMath ID 20001087 · AIM-COMBINATORICS-0212 · Combinatorics · AIM Workshop Problem Lists}
\subsection{Definitions and mathematical statement}
For a finite set $A\subseteq\mathbb Z$, its sumset is
$A+A=\{a+b:a,b\in A\}$; equal summands are permitted and each
resulting integer is counted once. Say $A$ is progression-free if
there are no three distinct elements $x,y,z\in A$ with $x+z=2y$.
Equivalently, $a,a+d,a+2d$ cannot all lie in $A$ for any integer
$d\ne0$. For each positive integer $n$, define
\[
 f(n)=\min\{|A+A|:A\subseteq\mathbb Z, |A|=n,
                                   A\text{ is progression-free}\}.
\]
The family is nonempty, and possible sumset sizes are positive
integers, so the minimum is well defined even though the ambient
integers are unbounded.

Freiman's problem asks how small $f(n)$ can be, especially its sharp
growth as $n\to\infty$. One may equivalently seek the smallest
possible doubling ratio $f(n)/n$ under the progression-free condition.
A lower bound must hold for every eligible integer set; a matching
upper bound requires actual sets of size $n$ with that small a sumset.
No upper bound on the diameter or largest element of $A$ is part
of the hypothesis. Translating or multiplying every element by the
same nonzero integer preserves both the condition and the sumset
cardinality. The target therefore concerns the additive structure of
the set rather than how it has been placed in an interval.
\subsection{Short English statement}
What is the smallest sumset of an n-element integer set containing no nonconstant three-term arithmetic progression?
\subsection{Sources}
\begin{itemize}
\item[S1] ULAM.AI and the UnsolvedMath Contributors, supplied problems.json, record 20001087 (AIM-COMBINATORICS-0212), AIM Workshop Problem Lists. Catalogue identity and source transcription; CC BY 4.0.
\url{https://huggingface.co/datasets/ulamai/UnsolvedMath}.
\item[S2] American Institute of Mathematics, Recent trends in additive combinatorics, item 1.6. Original workshop question underlying AIM-COMBINATORICS-0212.
\url{https://aimath.org/WWN/additivecomb/additivecomb.pdf}.
\end{itemize}
\end{document}
